\chapter{Project Overview}\label{ch:overview}

This chapter sets out what the project is and why it is worth doing. It defines
the research questions you will answer and the plan for answering them. Read it
first, and come back to it whenever you lose sight of why a step matters.

% =========================================================================
\section{Where you are starting from}\label{sec:overview:start}

Your thesis~\cite{burrell2026thesis} built a real-time speech enhancer for
hearing aids that is controlled by plain language. It runs two pipelines
concurrently:

\begin{itemize}
  \item \textbf{An audio pipeline.} It enhances the microphone signal frame by
  frame: a sliding STFT with $\sqrt{\text{Hann}}$ windows and weighted
  overlap-add, IMCRA noise estimation~\cite{cohen2003imcra}, and a parametric
  Wiener filter $G=(\xi/(1+\xi))^\beta$ with a decision-directed a priori
  SNR~\cite{ephraim1984mmse}.
  \item \textbf{An NLP pipeline.} It turns a typed complaint into an intent
  class and a severity score. It uses \code{all-mpnet-base-v2} sentence
  embeddings~\cite{reimers2019sbert,song2020mpnet}, a linear SVM (test
  macro-F1 0.957) and a Ridge regressor (CV MAE 0.100).
\end{itemize}
A parameter controller maps the result to DSP updates that take effect without
interrupting the audio.

The thesis answered its question: natural language \emph{can} drive the DSP, and
the path from sentence to sound is intact. It also listed limitations
honestly. One of them is the starting point for this project:

\begin{quote}\itshape
\code{all-mpnet-base-v2} is 110M parameters and effectively wants a GPU --- fine
for a laptop proof-of-concept, far beyond embedded hearing-aid hardware.
\end{quote}

The same is true of the audio pipeline, for less obvious reasons. It runs in
64-bit floating point, keeps about \num{23600} words of state between frames,
performs 12 divisions and two transcendental functions per frequency bin, and
delays every sample by \qty{21.33}{\ms} (\cref{ch:workload}). None of that
matters on a laptop. All of it matters on a device that has to run all day on
a hearing-aid battery, where the signal processor gets a power budget measured
in milliwatts or less.

% =========================================================================
\section{The core idea}\label{sec:overview:idea}

\Cref{ch:workload} shows that your system already has the structure of a real
hardware system. A \emph{data plane} (the audio path) runs every
\qty{10.667}{\ms} under a hard deadline. A \emph{control plane} (the intent path)
runs only when the user speaks. The interface between them is tiny: four
numbers ($\beta$, gain floor, output gain, tilt).

In hardware that interface becomes even smaller. Because $\beta$ and the floor
change only when the user complains, the whole gain curve can be precomputed by
the control plane into a small lookup table. The datapath then never evaluates
\code{pow} at all. Tilt and output gain become two registers. The NLP side
controls the DSP by \emph{writing a table and two registers}, which is exactly
how fitting parameters are programmed into real hearing-aid chips.

That gives the hardware version of your thesis question:

\begin{rqbox}{The project question}
What is the smallest, lowest-power hardware that can still do intent-controlled
speech enhancement, and what does each design decision cost in quality, latency,
memory, area and energy?
\end{rqbox}

% =========================================================================
\section{Research questions and hypotheses}\label{sec:overview:rq}

The project question breaks down into four research questions. Each one has a
measurable target, so you will know when you have answered it.

\begin{rqbox}{RQ1 --- Precision}
What minimum arithmetic precision (word lengths, number formats, function
approximations) keeps the IMCRA + parametric-Wiener pipeline perceptually
equivalent to the floating-point thesis reference?

\textbf{Target:} $|\Delta\text{PESQ}|\le0.1$ and $|\Delta\text{STOI}|\le0.01$
against float on the VoiceBank+DEMAND test set~\cite{valentini2016voicebank}.
\textbf{Answered in} \cref{ch:codesign} (Phase~1).
\end{rqbox}

\begin{rqbox}{RQ2 --- Latency}
How far can algorithmic latency be cut from \qty{21.33}{\ms} toward hearing-aid
territory (\qty{10}{\ms} or less, ideally about \qty{8}{\ms}), and what does each
step cost in quality and memory?

\textbf{Answered in} \cref{ch:codesign} (Phase~1), with hardware cost confirmed
in \cref{ch:rtl,ch:fpga}.
\end{rqbox}

\begin{rqbox}{RQ3 --- Cost}
What are the area, power, energy per frame and memory costs of the resulting
datapath on an FPGA and in an open-source ASIC flow? How do they trade against
clock frequency and parallelism?

\textbf{Answered in} \cref{ch:rtl,ch:fpga,ch:asic} (Phases~2, 4 and~5B).
\end{rqbox}

\begin{rqbox}{RQ4 --- The intent path at the edge}
How small and how low-precision can the intent path be while keeping macro-F1
close to the thesis's 0.957? The options are a quantised classifier head, a
smaller encoder, or an intent model that works directly on audio.

\textbf{Target:} within 0.02 macro-F1 for the quantised head; for encoders,
report the full trade-off curve.
\textbf{Answered in} \cref{ch:codesign} (first result) and \cref{ch:mlaccel}
(Phase~5A).
\end{rqbox}

The research questions come with four hypotheses that give the project its
first concrete experiments (\cref{tab:overview:hyp}). A hypothesis is a
prediction you could be wrong about. If one fails, you have still learned
something you can report.

\begin{table}[htbp]
\centering\small
\caption{Hypotheses, and how each is tested.}
\label{tab:overview:hyp}
\begin{tabularx}{\textwidth}{@{}lLLl@{}}
\toprule
 & Hypothesis & Test & Where \\
\midrule
H1 & A datapath of 16 bits or fewer, with LUT-based nonlinear functions, meets the RQ1 target. & Word-length sweep of the bit-accurate model & \cref{ch:codesign} \\
H2 & Removing the 2$\times$ zero padding changes PESQ by less than 0.05, while halving FFT work and per-bin state. & Frame-configuration study & \cref{ch:codesign} \\
H3 & Because the six classes are linearly separable, an int8 classifier head loses less than 0.01 macro-F1. & Post-training quantisation of the heads & \cref{ch:codesign,ch:mlaccel} \\
H4 & With resource sharing (one butterfly, one per-bin pipeline), the design meets the real-time deadline at a clock of \qty{5}{\MHz} or less. & Cycle count in simulation, then on the FPGA & \cref{ch:rtl,ch:fpga} \\
\bottomrule
\end{tabularx}
\end{table}

% =========================================================================
\section{Scope and non-goals}\label{sec:overview:scope}

A research project needs edges. These are deliberately outside it:

\begin{itemize}
  \item \textbf{The perceptual calibration study.} The thesis left the mapping
  from severity to parameter step size uncalibrated, and it needs a listening
  study with hearing-aid users. This project does not attempt that. It does make
  that study cheaper later: once parameter updates are just table writes, a
  calibration experiment can sweep them on real hardware.
  \item \textbf{Speech recognition.} The control plane keeps typed text, or
  replaces it with a small audio-domain intent model (\cref{ch:mlaccel}). It does
  not include a full ASR system.
  \item \textbf{Analog front end and acoustics.} There is no microphone
  preamplifier, ADC design, or feedback cancellation.
  \item \textbf{A product.} The aim is measured evidence and transferable
  skills, not a certified medical device.
\end{itemize}

% =========================================================================
\section{Deliverables}\label{sec:overview:deliverables}

\begin{enumerate}
  \item \textbf{A public repository}, \file{intent-se-hw}, containing:
        a bit-accurate Python golden model; SystemVerilog RTL; cocotb
        testbenches; FPGA project scripts; and, on the ASIC track, LibreLane
        configuration and reports.
  \item \textbf{A specification} (\code{spec-v1}): number formats, frame
        configuration and approximation choices, each backed by an experiment.
  \item \textbf{A real-time FPGA demonstration} on a PYNQ-Z2. Typed complaints
        change the audio you hear, through hardware registers. Record a short
        video.
  \item \textbf{A results set}: quality versus precision; latency versus quality
        and memory; FPGA resources and power; and either ML-head accelerator
        numbers or ASIC area, timing and power.
  \item \textbf{A written report} in paper format (\cref{ch:writeup}), suitable
        for a student conference submission or as the basis of an MSc thesis
        proposal.
\end{enumerate}

% =========================================================================
\section{Roadmap at a glance}\label{sec:overview:roadmap}

The plan assumes about 12--15 hours per week for 26 weeks, including two buffer
weeks (\cref{fig:overview:roadmap,tab:overview:phases}). \Cref{ch:plan} turns it
into a week-by-week schedule with a risk register, and shows how to shrink it if
you have less time.

\begin{figure}[htbp]
\centering
\begin{tikzpicture}[x=5.25mm,y=6.2mm,font=\sffamily\scriptsize]
  % week grid
  \foreach \w in {1,...,26}{
    \node[text=conceptc] at (\w-0.5,0.55) {\w};
    \draw[conceptc!25] (\w-1,0.2) -- (\w-1,-8.2);
  }
  \draw[conceptc!25] (26,0.2) -- (26,-8.2);
  \node[anchor=east,text=conceptc] at (-0.2,0.55) {week};
  % bars: name / start / end / row / colour
  \foreach \n/\s/\e/\r/\c in {
      {P0 Setup}/1/2/1/brand,
      {P1 Co-design + golden model}/3/7/2/brand,
      {P2 RTL design}/6/12/3/stepc,
      {P3 Verification}/7/12/4/stepc,
      {P4 FPGA prototype}/12/16/5/whyc,
      {P5 Track A (ML) or B (ASIC)}/17/22/6/rqc,
      {Write-up}/21/24/7/readc,
      {Buffer}/25/26/8/conceptc}{
    \fill[\c!75,rounded corners=0.6mm] (\s-1,-\r+0.32) rectangle (\e,-\r-0.32);
    \node[anchor=west,text=black] at (\e+0.15,-\r) {\n};
  }
  % milestones
  \foreach \m/\w in {M0/2,M1/7,M2/12,M3/16,M4/22,M5/24}{
    \draw[pitc,line width=0.6pt,dashed] (\w,0.2) -- (\w,-8.4);
    \node[fill=pitc,text=white,rounded corners=0.5mm,inner sep=1.2pt,font=\sffamily\bfseries\tiny] at (\w,-8.8) {\m};
  }
\end{tikzpicture}
\caption{The project at a glance: phases (bars), weeks, and milestones M0--M5
(red, at the end of the marked week).}
\label{fig:overview:roadmap}
\end{figure}

\begin{table}[htbp]
\centering\small
\caption{Phases, chapters and milestones.}
\label{tab:overview:phases}
\begin{tabularx}{\textwidth}{@{}llllL@{}}
\toprule
Phase & Chapter & Weeks & Milestone & Key deliverable \\
\midrule
0 Setup & \cref{ch:setup} & 1--2 & M0 & Toolchain, data, float baseline metrics \\
1 Co-design & \cref{ch:codesign} & 3--7 & M1 & Bit-accurate model, design-space results, golden vectors, \code{spec-v1} \\
2 RTL & \cref{ch:rtl} & 6--12 & \multirow{2}{*}{M2} & \multirow{2}{=}{RTL bit-exact against the golden model in simulation} \\
3 Verification & \cref{ch:verification} & 7--12 & & \\
4 FPGA & \cref{ch:fpga} & 12--16 & M3 & Real-time demo; latency, resource and power measurements \\
5A ML accel. & \cref{ch:mlaccel} & 17--22 & M4 & Int8 head engine plus encoder or audio-intent study \\
5B ASIC & \cref{ch:asic} & 17--22 & M4 & Area, timing and power from an open PDK flow \\
Write-up & \cref{ch:writeup} & 21--24 & M5 & Report, release v1.0, demo video \\
\bottomrule
\end{tabularx}
\end{table}

% =========================================================================
\section{Two tracks after the prototype}\label{sec:overview:tracks}

After the FPGA demo (M3) you choose one of two directions for six weeks. Choose
around week 16, based on what you enjoyed and what you want to do next:

\begin{description}[style=nextline,leftmargin=1.2em]
  \item[Track A --- ML acceleration (\cref{ch:mlaccel})]
  Bring the intent path into hardware: an int8 matrix--vector engine for the
  classifier and severity heads, then either a compressed text encoder or a
  small model that recognises spoken intent directly from audio features
  computed by your own FFT. Choose this if you are aiming at AI-accelerator or
  edge-ML roles.
  \item[Track B --- VLSI and ASIC (\cref{ch:asic})]
  Take the per-bin datapath through an open-source RTL-to-GDS flow and report
  area, timing and power in a real process. Optionally submit a small block to
  Tiny Tapeout. Choose this if you are aiming at digital IC design or DSP-ASIC
  roles.
\end{description}

Both tracks reuse everything from Phases 0--4, and both answer a research
question (RQ4 or RQ3) with numbers.

% =========================================================================
\section{Skills you will gain}\label{sec:overview:skills}

\begin{table}[htbp]
\centering\small
\caption{Skills, where you practise them, and the evidence you can show.}
\label{tab:overview:skills}
\begin{tabularx}{\textwidth}{@{}lLL@{}}
\toprule
Skill & Where in the project & Evidence for employers \\
\midrule
Fixed-point and word-length design & Phase~1 golden model and sweeps & Quality-versus-bits plots; a justified spec \\
Algorithm--architecture co-design & Phase~1 (LUT gain, log-domain trackers, frame study) & \enquote{Decisions that flip} story; Pareto plots \\
Microarchitecture and RTL (SystemVerilog) & Phase~2 & Lint-clean RTL; resource log; block diagrams \\
Verification (cocotb, Verilator, assertions) & Phase~3 & Bit-exact regression in CI; coverage report \\
FPGA prototyping (Zynq, AXI, I2S) & Phase~4 & Working demo video; timing and utilisation reports \\
ML hardware (quantisation, MAC arrays) & Track A & Int8 accuracy study; cycles and energy per inference \\
Physical design (synthesis, place and route, STA, power) & Track B & GDS layout; area, timing and power reports \\
Research method and writing & Throughout & Lab notebook, reproducible results, paper draft \\
\bottomrule
\end{tabularx}
\end{table}

% =========================================================================
\section{Why this matters beyond the project}\label{sec:overview:industry}

Hearing aids are one of the most demanding applications of low-power digital
signal processing. The devices are tiny, run all day on a small battery, and
have to process audio with very little delay. Hearing-aid manufacturers
therefore design their own low-power DSP chips. Neural networks now run on the
device itself: Oticon, for example, marketed Oticon More (2021) as the first
hearing aid with an on-board deep neural network. Denmark is home to several of
the large hearing-aid companies, including Demant (Oticon), GN and WS
Audiology.

Your thesis was done with Oticon. A project that goes from that thesis to
measured, verified hardware gives you a concrete story for exactly the
engineering teams that build these chips: fixed-point DSP, low-power
architecture, verification, and ML at the edge.

\begin{donebox}[Before moving on]
\begin{checklist}
  \item You can state the project question and RQ1--RQ4 in your own words.
  \item You know which chapter answers each research question.
  \item You have read \cref{ch:workload} (the measured workload) and skimmed
        \cref{ch:plan} (the schedule).
  \item You have discussed scope with your supervisor, especially Track A
        versus Track B and the credit load (\cref{ch:plan}).
\end{checklist}
\end{donebox}
